\section{Berücksichtigung der Normalkraft und der Bohrung}
\subsection{Fläche und FTM des Profils}
Die Bohrung wird Material in dem gefährlichem Querschnitt abtrennen. Die Bohrung ist symmetrisch um die y-Achse, und geht durch eine Trapezförmige Querschnitt des Profils.

\begin{figure}[H]
  \centering
  \begin{tikzpicture}
    \draw[blue] (2.5, 2) -- (2.5, 4);
    \fill[pattern=horizontal lines, pattern color=blue] (2.5, 2) -- (3, 4) -- (2.5, 4);

    \draw[blue] (7.5, 2) -- (7, 0);
    \fill[pattern=horizontal lines, pattern color=blue] (7.5, 2) -- (7, 0) -- (7.5, 0);

    \draw[ultra thick] (0,0) -- (2,0) -- (3, 4) -- (0, 4) -- (0,0);
    \draw[ultra thick] (5,0) -- (7.5,0) -- (7.5, 4) -- (5, 4) -- (5,0);
  \end{tikzpicture}
  \caption{Vereinfachung des Profils}
  \label{fig:trapezinrechteck}
\end{figure}
Wenn die Biegung um die z-Achse passiert, kann das Querschnitt als im \figref{fig:trapezinrechteck} zu sehen ist vereinfacht werden. Dieser Vereinfachung ändert $I_z$ nicht.

Das Flächenträgheitmoment den gebohrten Querschnitte nach dem Formel:
\begin{align*}
  &I_{bohr} = 2 \cdot \frac{d^3 \cdot t}{12} = 138666,663 \, mm^4
\end{align*}

Die neue Fläche muss auch berechnet werden:
\begin{align*}
  &A = A_{U240} - A_{bohr} = 4230 - 2 \cdot 40 \cdot 13 = 3190 \, mm^2 \\
\end{align*}

Der neue Schwerpunkt kann von dem altem Schwerpunkt und der Schwerpunkt der Bohrung berechnet werden.
\begin{align*}
  s_z &= \frac{A_{U240} \cdot e_y - A_{bohr} \cdot \frac{b}{2}}{A} \\
  &= \frac{4230 \cdot 22{,}3 - 1040 \cdot 42{,}5}{3190} \\
  &= 15{,}71 \, \text{mm}
\end{align*}

Die Flächenträgheitmomente kann in diesem Schwerpunkt verschoben werden mit hilfe der Steiner-Satz.
\begin{align*}
  I_{S,bohr} &= I_{bohr} + A_{bohr} \cdot (\frac{b}{2} - s_z)^2 \\
  &= 885078{,}93 \, \text{mm}^4\\
  \\
  I_{S,U240} &= I_{U240} + A_{U240} \cdot (e_y - s_z)^2 \\
  &= 2480000 + 4230 \cdot (22{,}3 - 15{,}71)^2 \\
  &= 2663700{,}86 \, mm^4 \\
\end{align*}

Das Flächenträgheitmoment des neues Querschnitts kann jetzt berechnet werden.
\begin{align*}
  I_{z} &= I_{S,U240} - I_{S,bohr} \\
  I_{z} &= 2663700{,}86 - 885078{,}93 \\
  I_{z} &= 1778621{,}93 \, \text{mm}^4
\end{align*}

\subsection{Spannungen in dem Querschnitt}
Die Normalspannung wegen die Normalkraft:
\begin{align*}
  \sigma_{x,N} = \frac{N}{A} = \frac{10464 \, N}{3190 \, mm^2} = 3,28 \, \text{MPa}
\end{align*}

Die Normalspannung wegen der Biegemoment:
\begin{align*}
  \sigma_{x,b}(z) &= \frac{M_{b}^{(1)}(a)}{I_{z}} \cdot z \\
  &= \frac{7 \cdot 10^6 \, \text{Nmm}}{1778621{,}93 \, \text{mm}^4} \cdot z \\
  &= 3{,}936 \cdot z \, \text{[MPa]}
\end{align*}

Nach dem Superpositionsprinzip können die Zwei Spannungswerte addiert, so bekommt man die Spannunsverteilung in diesem Querschnitt:
\begin{align*}
  \sigma_{x}(z) &= \sigma_{x,N} + \sigma_{x,b}(z) \\
  &= 3{,}936z + 3,28 \, \text{[MPa]}
\end{align*}

\subsection{Dimensionierung}
Mit dem neuen Schwerpunkt ($s_z = 15,71$mm) ist die neue Randfase $e = -69,29$mm

Mit Substition bekommt die maximale Spannung in dem Querschnitt.
\begin{align*}
  \sigma_{x}(-69,29) &= -269,45 \, \text{MPa}
\end{align*}

Dieser Wert is deutlich größer, als die zulassige Spannung $\sigma_{zul} = 177,5$ MPa. Man muss also ein neues Profil ausgewählt werden.
